\documentclass[12pt]{article}

\usepackage[utf8]{inputenc}
\usepackage{amsmath, amssymb, amsthm}
\usepackage{geometry}
\usepackage[T1]{fontenc}
\usepackage{lmodern}

\geometry{a4paper, margin=2.5cm}

\title{Factori Invarianti, Divizori Elementari și Module Indecompozabile}
\author{}
\date{}

\theoremstyle{definition}
\newtheorem{definition}{Definiție}

\theoremstyle{plain}
\newtheorem{theorem}{Teoremă}

\theoremstyle{remark}
\newtheorem{remark}{Observație}

\theoremstyle{plain}
\newtheorem{proposition}{Propoziție}

\begin{document}

\maketitle

\section*{Factori invarianti}

\begin{definition}
Fie $R$ un PID și $M$ un $R$-modul finit generat. O \emph{reprezentare în factori invarianti} este o scriere


\[
M \cong R^k \oplus R/d_1R \oplus \cdots \oplus R/d_mR,
\]


unde $d_1, \dots, d_m \in R \setminus \{0\}$ satisfac


\[
d_1 \mid d_2 \mid \cdots \mid d_m.
\]


Elementele $d_1, \dots, d_m$ se numesc \emph{factori invarianti}.
\end{definition}

\begin{remark}
Condiția de divizibilitate $d_1 \mid d_2 \mid \cdots \mid d_m$ face reprezentarea \emph{unică} până la unități.
\end{remark}

\begin{proposition}
Factorii invarianti descriu modulul torsional ca sumă directă ordonată:


\[
t(M) \cong R/d_1R \oplus \cdots \oplus R/d_mR.
\]


\end{proposition}

\section*{Divizori elementari}

\begin{definition}
Fie $R$ un PID. O \emph{reprezentare în divizori elementari} este o scriere


\[
M \cong R^k \oplus R/p_1^{e_1}R \oplus \cdots \oplus R/p_s^{e_s}R,
\]


unde $p_i$ sunt primi (până la unități), iar $e_i \ge 1$.
\end{definition}

\begin{remark}
Divizorii elementari sunt unici până la permutare și unități. Ei corespund componentelor $p$-primare.
\end{remark}

\begin{theorem}[Legătura dintre factori invarianti și divizori elementari]
Dacă


\[
d_i = p_{i1}^{e_{i1}} \cdots p_{ik_i}^{e_{ik_i}}
\]


este factorizarea în factori primi, atunci


\[
R/d_iR \cong R/p_{i1}^{e_{i1}}R \oplus \cdots \oplus R/p_{ik_i}^{e_{ik_i}}R.
\]


Astfel, factori invarianti determină divizorii elementari și invers.
\end{theorem}

\begin{remark}
Reprezentarea în factori invarianti este „globală”, iar cea în divizori elementari este „locală” (pe fiecare prim).
\end{remark}

\section*{Module decompozabile și indecompozabile}

\begin{definition}
Un modul $M$ se numește \emph{decompozabil} dacă există submodule nenule $M_1, M_2$ cu


\[
M = M_1 \oplus M_2.
\]


Dacă nu există o astfel de descompunere, modulul se numește \emph{indecompozabil}.
\end{definition}

\begin{proposition}
Dacă $R$ este un PID, atunci modulele indecompozabile finit generate sunt exact cele de forma


\[
R/p^eR,
\]


unde $p$ este prim și $e \ge 1$.
\end{proposition}

\begin{remark}
Modulele de forma $R/p^eR$ sunt „blocurile atomice” din care se construiesc toate modulele finit generate.
\end{remark}

\begin{theorem}[Decompoziția în indecompozabile]
Fie $R$ un PID și $M$ un $R$-modul finit generat. Atunci


\[
M \cong R^k \oplus \bigoplus_{i=1}^s R/p_i^{e_i}R,
\]


iar această descompunere este unică până la permutarea factorilor indecompozabili.
\end{theorem}

\begin{remark}
Aceasta este versiunea pentru PID a teoremei Krull--Schmidt.
\end{remark}

\end{document}
